\begin{titlepage}
  \setlength{\parindent}{0cm}

  \includegraphics[width=200px]{Images/uob-logo-black-transparent}

  \Large
  Faculty of Engineering and Design

  \vspace{20pt}

  \begin{singlespace}
  \LARGE
  Final Year MEng Project \\
  Project Report
  \end{singlespace}

  \vspace{40pt}
  \textbf{Robot Football - High Level Control and Strategy Implementation}

  \vspace{10pt}
  \emph{Anthony Richards} \\
  \emph{\today}
  
  \vspace{20pt}
  \begin{singlespace}
  \emph{ \large
  `I certify that I have read and understood the entry in the Student Handbook
  for the Department of Mechanical Engineering on Cheating and Plagiarism and
  that all material in this assignment is my own work, except where I have
  indicated with appropriate references. I agree that, in line with Regulation
  15.3(e), if requested I will submit an electronic copy of this work for
  submission to a Plagiarism Detection Service for quality assurance purposes.'
 }
  \end{singlespace}
  
  \vspace{20pt}
  
  {\large
  Author's signature
  \ldots\ldots\ldots\ldots\ldots\ldots\ldots\ldots\ldots\ldots\ldots\ldots\ldots
  \ldots\ldots
  }

  \vspace{80pt}
  Supervisor: \emph{Dr.~Pejman Iravani}

  \vspace{10pt}
  Assessor: \emph{Dr.~Francis Robinson}
\end{titlepage}

\pagenumbering{roman}

\begin{abstract}

The aim of this project is to develop of a set of control algorithms for the
Federation of International Robot-soccer Association (FIRA) Middle League
Simursot Game, a simulated five-a-side game played by small (not humanoid)
autonomous robots.  Working from an idea proposed by one of the papers read in a
previous report, this design project has produced a system that operates using
potential field force navigation.  This method uses sets of attractive and
repulsive fields to push and pull the robots around the field, in order to
achieve the game's goals.  The system described scores a goal under a given set
of circumstances, and provides a base for a more complex system that can be
developed to play a game of football against an opposing team.

\end{abstract}

\tableofcontents

\listoffigures

\section*{Document Conventions}
The following unit conventions are used by this document:

\begin{itemize}
  \item Unless otherwise stated, all units used are standard SI units
  \item The symbol `\si{\inch}' is used to designate the unit inches
\end{itemize}

All images displaying field layouts are shown as a greyscale image representing
relative magnitude.  A white pixel represents a point of highest magnitude
(either gradient or field intensity), while a black pixel represents a point of
lowest magnitude.

\subsection*{List of Acronyms}

The following acronyms are in use in this document:

\begin{itemize}
  \item CPU - Central Processing Unit
  \item CRT - C Runtime Library
  \item CSV - Comma Separated Variable
  \item DLL - Dynamic Link Library
  \item FIRA - Federation of International Robot-Soccer Association
  \item GPR - General Purpose Register
  \item GPU - Graphics Processing Unit
  \item SIMD - Single Instruction Multiple Data
\end{itemize}

\subsection*{List of Symbols}

The following symbols are in use in this document:

\begin{itemize}
  \item $\boldsymbol{F}(x,y)$ - The gradient of a potential field at point $\left(x,y\right)$.  The real part is the gradient in $x$, while the imaginary part is the gradient in $y$.
  \item $G(s)$ - A transfer function, represented in the Laplace domain.
  \item $k_{attr}$ - A weight for an attractive field.
  \item $k_{repulse}$ - A weight for a repulsive field.
  \item $N$ - The number of elements in a summation.
  \item $p_i(x,y)$ - The intensity of potential field $i$ at point $\left(x,y\right)$
  \item $P(x,y)$ - The total potential field intensity at point $\left(x,y\right)$.
  \item $R_0$ - The limiting radius of a field.
  \item $R^2$ - A measure of the quality of a line of best fit, where 1 is a perfect fit a smaller number represents a worse fit.
  \item $t$ - The time in seconds (normally relative to the beginning of the operation).
  \item $\sigma$ - The variable controlling the spread of a Gaussian function
\end{itemize}

\cleardoublepage

\pagenumbering{arabic}